\section{Keep the six calls and their result units explicit}
\texttt{wav\_parse} takes a byte pointer in \texttt{a0}, an unsigned complete byte count
in \texttt{a1} and a descriptor destination in \texttt{a2}. It returns the frame count
in \texttt{v1} after successful publication. \texttt{wav\_bound} takes a pointer to the
sixteen-byte description prefix in \texttt{a0}; its successful \texttt{v1} is the
canonical file size, $44+2CN$. It validates rate, channels and frame count before using
them in that expression.

\texttt{wav\_filter} takes interleaved PCM, a description prefix, a disjoint PCM
output and eight signed Q15 coefficient codes in \texttt{a0} through \texttt{a3}.
Its successful \texttt{v1} counts clipped sample results across all channels.
\texttt{wav\_encode} takes PCM, a description prefix, a disjoint byte destination
and unsigned byte capacity in the same four argument registers. Its successful
\texttt{v1} is the complete encoded file length.

\texttt{wav\_load} takes a path, byte buffer, unsigned buffer capacity and descriptor
in \texttt{a0} through \texttt{a3}; \texttt{wav\_save} takes a path, complete bytes
and their unsigned count in \texttt{a0} through \texttt{a2}. Both return complete file
bytes in \texttt{v1}. None of those byte counts can be substituted for a filter frame
count or a clipped-sample count merely because the values share a result register.

Every completed call returns status zero in \texttt{v0} on success. Format or
unsupported-description failures return one; excessive complete byte counts or invalid
file/output capacities return two; file I/O failures return three. Every failure clears
\texttt{v1}. A valid zero-frame parse and a successful filter with no clips both have
zero \texttt{v1}; inspect status and the call's unit before interpreting that zero.
